% Remove any automatic bibliography title and formatting
\makeatletter
\renewenvironment{thebibliography}[1]
     {\list{\@biblabel{\@arabic\c@enumiv}}%
           {\settowidth\labelwidth{\@biblabel{#1}}%
            \leftmargin\labelwidth
            \advance\leftmargin\labelsep
            \@openbib@code
            \usecounter{enumiv}%
            \let\p@enumiv\@empty
            \renewcommand\theenumiv{\@arabic\c@enumiv}}%
      \sloppy
      \clubpenalty4000
      \@clubpenalty \clubpenalty
      \widowpenalty4000%
      \sfcode`\.\@m}
     {\def\@noitemerr
       {\@latex@warning{Empty `thebibliography' environment}}%
      \endlist}
\makeatother
\begin{thebibliography}{99}
% IEEE-style references, numbered to match in-body [N] citations.

\bibitem{qi2019}
Q. Qi, F. Tao, T. Zuo and A. Y. Nee, ``A review of digital twins in smart industries,'' \textit{IEEE Access}, vol. 7, pp. 14232--14252, Jan. 2019.

\bibitem{fuller2020}
A. Fuller, Z. Fan, C. Day and C. Barlow, ``Digital Twin: Enabling Technologies, Challenges and Open Research,'' \textit{IEEE Access}, vol. 8, pp. 108952--108971, June 2020.

\bibitem{geyer2021}
F. Geyer and S. Bondorf, ``Proximity-based service discovery for distributed digital twin systems,'' \textit{IEEE Access}, vol. 9, pp. 12456--12470, 2021.

\bibitem{lv2022}
Z. Lv, D. Chen, R. Lou, and H. Song, ``Digital Twins: A Survey on Enabling Technologies, Challenges, Trends, and Future Prospects,'' \textit{IEEE Communications Surveys \& Tutorials}, vol. 24, no. 1, pp. 54--87, 2022.

\bibitem{grieves2022}
M. Grieves and J. Vickers, ``Digital twins: Recent advances and future directions in engineering fields,'' \textit{Nature Reviews Methods Primers}, vol. 2, no. 1, pp. 1--18, Dec. 2022.

\bibitem{wang2022}
J. Wang, L. Ye and R. X. Gao, ``A simulation-based digital twin model for data-driven decision optimization,'' \textit{Journal of Manufacturing Systems}, vol. 62, pp. 245--258, Jan. 2022.

\bibitem{kim2022}
S. S. Kim, Y. J. Lee and M. G. Lee, ``Digital twin model with machine learning for real-time monitoring and predictive maintenance,'' \textit{International Journal of Precision Engineering and Manufacturing}, vol. 23, no. 5, pp. 565--578, May 2022.

\bibitem{xu2023}
C. Xu, Z. Tang, H. Yu, P. Zeng and L. Kong, ``Digital Twin-Driven Collaborative Scheduling for Heterogeneous Task and Edge-End Resource via Multi-Agent Deep Reinforcement Learning,'' \textit{IEEE Journal on Selected Areas in Communications}, vol. 41, no. 10, pp. 3056--3069, Oct. 2023.

\bibitem{dai2024}
Y. Dai, J. Zhao, J. Zhang, Y. Zhang and T. Jiang, ``Federated Deep Reinforcement Learning for Task Offloading in Digital Twin Edge Networks,'' \textit{IEEE Transactions on Network Science and Engineering}, vol. 11, no. 3, pp. 2849--2861, May-June 2024.

\bibitem{sai2024}
A. M. V. V. Sai, C. Wang, Z. Cai and Y. Li, ``Navigating the Digital Twin Network landscape: A survey on architecture, applications, privacy and security,'' \textit{High-Confidence Computing}, vol. 4, no. 4, p. 100269, Dec. 2024.

\bibitem{raza2025}
S. M. Raza, R. Minerva, N. Crespi, M. Alvi, M. Herath and H. Dutta, ``A comprehensive survey of Network Digital Twin architecture, capabilities, challenges, and requirements for Edge-Cloud Continuum,'' \textit{Computer Communications}, vol. 236, pp. 165--192, Jan. 2025.

\bibitem{trandang2025}
H. Tran-Dang and D.-S. Kim, ``Digital Twin-empowered intelligent computation offloading for edge computing in the era of 5G and beyond: A state-of-the-art survey,'' \textit{ICT Express}, vol. 11, no. 2, pp. 167--180, Apr. 2025.

\bibitem{esen2026}
E. Esen, A. Akbulut and C. Catal, ``Chaos experiments in microservice architectures: A systematic literature review,'' \textit{Computer Standards \& Interfaces}, vol. 97, p. 104116, 2026.

\end{thebibliography}
